\section{Additional operator definitions}
\subsection{Architecture and workflow}
The code separates lattice/grid geometry, MPI exchange, parameter and field
I/O, and fluid collision operators. A byte-valued voxel file assigns initial
phases and solid or boundary labels. A parameter file selects the lattice,
relaxation operator, relaxation times, forcing and wetting. D2Q9 and D3Q19,
SRT and raw-moment MRT, and the single-component and color-fluid executables
are exercised here as configurations of the same executable. The Shan--Chen
placeholder is rejected explicitly.

Directional populations reside in arrays; streaming exchanges the required
neighbor data between MPI subdomains. Macroscopic fields are written as VTI
pieces that the analysis reconstructs into a complete global lattice.
Restart files contain outgoing populations and a sidecar recording the time,
lattice, phase count and next streaming action. Historical checkpoints
without an identifiable stage are rejected. A fresh convergence history is
required after restart. Exceptions and invalid populations cause a nonzero
MPI failure status.

The reproducibility scripts generate every geometry, record inputs and
executable hashes, preserve raw output and calculate independent reference
solutions. Each study uses a frozen executable so a later build cannot alter
an ongoing batch. Algebraic checks are available through CTest.
The public-rock preparation additionally pins the source DOI, data license,
checksums, crop coordinates, voxel scale and periodic connectivity. Numerical
timers include exchange and finite checks, while field/diagnostic output and
checkpoint I/O are counted separately.

\subsection{Bulk dynamics, forcing and viscosity}
In lattice units, the component densities and color field are
\begin{equation}
 \rho_k=\sum_i f_i^k,\quad \rho=\rho_0+\rho_1,\quad
 \chi=(\rho_0-\rho_1)/\rho,\quad
 \rho\vct u=\sum_i\vct e_i(f_i^0+f_i^1)+\vct F/2.
\end{equation}
Here $\vct F$ is force per unit volume. The mixture equilibrium is
\begin{equation}
 f_i^{eq}=\rho_0\phi_i^0+\rho_1\phi_i^1+
 \rho w_i[3\vct e_i\!\cdot\!\vct u+
 4.5(\vct e_i\!\cdot\!\vct u)^2-1.5|\vct u|^2].
\end{equation}
The phase weights give $c_{p,k}^2=3(1-\alpha_k)/5$ in D2Q9 and
$(1-\alpha_k)/2$ in D3Q19. The present quantitative tests use equal reference
densities and the standard pressure weights, so $p=\rho/3$.

The legacy arithmetic rule interpolates $\tau$ by component density fraction.
For $\varphi=\rho_0/\rho$ and $\nu_k=(\tau_k-1/2)/3$, the optional harmonic
kinematic-viscosity rule is
\begin{equation}
 \nu^{-1}=\varphi/\nu_0+(1-\varphi)/\nu_1,\qquad \tau=1/2+3\nu.
\end{equation}
The two rules are compared at equal density; equivalence to a dynamic-viscosity
law is not assumed at unequal density.

With $\vct m=M\vct f$, collision and body-force insertion become
\begin{equation}
 \vct m^*=\vct m-S(\vct m-\vct m^{eq})+(I-S/2)M\vct G,
 \quad G_i=w_i[3(\vct e_i-\vct u)+9(\vct e_i\!\cdot\!\vct u)\vct e_i]\!\cdot\!\vct F.
\end{equation}
This uses established discrete forcing and moment relaxation
\cite{guo2002,lallemand2000}. Conserved rates are zero. The retained
nonconserved rates are $s_u=1/\tau$ and $s_q=8(2-s_u)/(8-s_u)$, with
$(1/s_u-1/2)(1/s_q-1/2)=3/16$ \cite{ginzburg2003}.

Interfacial perturbation and recoloring retain the recovered structure:
\begin{align}
 \Delta_i&=\frac{9\sigma |\vct g|}{4\tau}
 \left[w_i(\vct e_i\!\cdot\!\vct n)^2-B_i\right],\quad
 \vct n=\vct g/|\vct g|,\\
 f_i^{0,*}&=\frac{\rho_0}{\rho}f_i^{post}+
 \beta\frac{\rho_0\rho_1}{\rho^2}\cos\varphi_i f_i^{eq}(\vct u=0),\qquad
 f_i^{1,*}=f_i^{post}-f_i^{0,*}.
\end{align}
The actual arrays satisfy zero perturbation mass and momentum and
$\sum_i\vct e_i\vct e_i\Delta_i=\sigma|\vct g|(\vct n\vct n-I)/(2\tau)$.
Every nonzero capillary moment is coupled to a mode with rate $s_u$.

Initialization constructs incoming rest populations with $-\vct G/2$ and
applies the normal collision kernel once. The stored state is then outgoing
at time zero. After streaming, a uniform forced domain has physical velocity
$\vct u(t)=t\vct F/\rho$. The same state convention is used for routine and
convergence-triggered checkpoints.

\subsection{Wetting and termination}
At a wall-adjacent fluid node, the historical gradient skips solid neighbors.
Writing $\mathcal F$ for the remaining stencil,
\begin{equation}
 \vct b=\sum_{i\in\mathcal F}w_i\vct e_i(\chi_i-\chi),\quad
 A=\sum_{i\in\mathcal F}w_i\vct e_i\vct e_i,
 \quad\vct g_{legacy}=3\vct b,\quad\vct g_{ren}=A^{-1}\vct b.
\end{equation}
The second expression is a standard weighted linear reconstruction, not a
new wetting model. It agrees with the bulk stencil and exactly reproduces
linear fields when $A$ has full rank. Unsupported rank-deficient stencils
are rejected. Both variants apply the same recovered angle rotation and
wall normal. On a flat lower wall, the legacy stencil scales the tangential
and normal derivatives by $5/6$ and $1/2$ respectively.

Steady termination requires finite state, stable requested scalar diagnostics,
velocity convergence and, for color flow, convergence of the spatial color
field, for three successive reports. Both report-to-report and adjacent-step
changes must pass; an even reporting interval cannot certify a period-two
oscillation as stationary. Constant saturation cannot terminate a
translating phase field. Periodic acceleration benchmarks disable steady
stopping; undefined pressure-based permeability is identified separately
from invalid numerical state.

\subsection{Experimental capillary operators}
Two optional controls examine limitations exposed by the static tests.
The first fits $\psi=\operatorname{atanh}(\chi)$ on available fluid neighbors
with the same moment tensor $A$, continues $\tanh(\psi+\vct e_i\cdot\nabla\psi)$
over a virtual full stencil, then applies the common angle rotation.
Color is clipped to $\pm(1-10^{-8})$ for this transform. Equilibrium-profile
continuation has prior work \cite{bhattacharya2026}; this local fit is not
an exact implementation of that paper's global ghost-node scheme.

The second replaces perturbation with established continuum surface forcing
\cite{lishchuk2003,sedahmed2024}:
\begin{equation}
 \vct F_s=\frac{\sigma}{2}\kappa\vct g,\qquad
 \kappa=-(I-\vct n\vct n):\nabla\vct n,\quad
 \vct n=\vct g/|\vct g|.
\end{equation}
Normal derivatives use the weighted fluid stencil and $A^{-1}$ at walls.
An additional MPI exchange supplies neighboring normals. Below
$|\vct g|=10^{-9}$, normal and capillary force are zero. The combined external
and capillary force enters both Guo collision and physical half-force
velocity; incoming initialization subtracts its capillary half-force before
the initial collision. Recoloring and the selected wall projection remain
common. This is a separate force-operator comparison, not a wall-only
ablation or replication of a published complete model.
The experimental force accepts closed VD/GV flow at equal reference
density and standard pressure weights. A guarded flat D2Q9 pressure
extension is assessed below: it requires ghost color, conservative phase
transfer and a pure force-free reservoir collar. Velocity reservoirs and
selective films remain unsupported. Neither option changes the legacy defaults.

An optional gradient form evaluates the same continuum force as
\begin{equation}
 \vct F_s=-\frac{\sigma}{2}\vct n
       [(I-\vct n\vct n):\nabla\vct g].
\end{equation}
It differentiates the unnormalized gradient with the same fluid stencil.
The identity follows from
$(I-\vct n\vct n):\nabla\vct g=|\vct g|(I-\vct n\vct n):\nabla\vct n$;
its discrete realization differs from normal differentiation.
Related Hessian forms and their discretization sensitivity are already
studied in phase-field LBM \cite{sugimoto2026}. This control introduces
neither an adjusted cutoff nor a fitted surface-tension coefficient.

\subsection{Pore-scale observables}
For $N$ total voxels, $N_f$ pore voxels, and $\varphi_0=\rho_0/\rho$,
$\varphi_1=1-\varphi_0$ in the fluid, the analysis defines
\begin{equation}
 S_k=\frac{1}{N_f}\sum_{fluid}\varphi_k,\qquad
 q_k=\frac{1}{N}\sum_{fluid}\varphi_k u_x,\qquad q=q_0+q_1.
\end{equation}
These are phase-fraction-weighted \emph{advective} Darcy fluxes of the common
mixture velocity. They do not include the component segregation momentum
introduced by recoloring. Component masses $\sum\rho_k$ are monitored
separately; in a weakly compressible closed domain, fixed mass does not make
the volume-fraction saturation exactly constant.

At unit reference density, the single-phase body-force estimate is
$K=\nu_{ref}q_{ref}/F_x$, and a phase normalization is
$k_{r,k}=\nu_kq_k/(F_xK)$. The latter is independently verified only for the
isolated pure-phase pore bundle below. The SCSP output $\vct u_{ref}$ uses
reference-density momentum and equilibrium, whereas color flow divides
momentum by local total density. A pure-phase comparison consequently uses
$\vct u_{ref}/\rho$ at unit reference density. The rock comparison retains
this convention and its small finite-density differences explicitly.

